\documentclass[11pt]{article}

\usepackage[T1]{fontenc}
\usepackage{amsmath,amssymb,amsthm}
\usepackage{booktabs}
\usepackage{hyperref}

\title{LRT-Based Modification Indices and Exact Expected Parameter Change}
\author{}
\date{\today}

\newcommand{\E}{\mathbb{E}}
\newcommand{\tr}{\operatorname{tr}}
\newcommand{\Var}{\operatorname{Var}}
\newcommand{\dto}{\xrightarrow{d}}
\newcommand{\pto}{\xrightarrow{p}}
\newcommand{\Th}{\hat\theta}

\theoremstyle{plain}
\newtheorem{theorem}{Theorem}
\newtheorem{proposition}{Proposition}
\theoremstyle{remark}
\newtheorem*{remark}{Remark}

\begin{document}
\maketitle

\section*{Purpose}

A modification index (MI) asks, for a single constrained parameter, how much fit
would improve if the constraint were freed. In measurement-invariance work the
constraint of interest is a cross-group equality (a loading or intercept tied
across groups), and the MI localizes \emph{which} item is non-invariant. magmaan
ships the classical one-step form: a Lagrange-multiplier (score) statistic
evaluated once at the constrained fit, with robust (sandwich-scaled) variants.

This note records the refit-based likelihood-ratio (LRT) counterpart, why it is
the correct tool when the model is structurally misspecified elsewhere, and what
it costs. The headline cost result is that the refit sweep is cheap: a full
invariance sweep over a 20-item two-group model runs in about one tenth of a
second, cold-started, in pure R.

\section{The three localizations and their anchor points}

Let $H_0$ (constrained) be nested in $H_1$ (the one extra constraint released),
with $m$ released degrees of freedom ($m=1$ for a single equality). The same
null can be probed from three places:
\begin{itemize}
\item \textbf{Score / LM (the MI).} Fit only $H_0$; test whether the score for
  the frozen direction is far from zero. One fit, anchored at the
  \emph{constrained} model.
\item \textbf{Wald.} Fit only $H_1$; test whether $R\Th_1 = q$. One fit, anchored
  at the \emph{free} model.
\item \textbf{LRT.} Fit \emph{both}; compare $T_d = T_0 - T_1 = 2N(F_0 - F_1)$.
\end{itemize}
All three converge to $\chi^2_m$ under the null and are asymptotically
equivalent. They differ in finite samples and, decisively, in their robustness
to misspecification \emph{elsewhere in the model}.

\section{Two senses of misspecification}

Write the candidate as direction $j$ and suppose the model is also wrong in some
other direction $k$.

\paragraph{Distributional misspecification} (non-normality, an inefficient
weight $W \neq \Gamma^{-1}$, the failure of the information equality
$\mathcal{I} = -\E[H]$). This inflates or deflates the \emph{reference
distribution} of the statistic. The robust score test repairs it by building the
meat from $\hat\Gamma$ rather than assuming $\mathcal{I}$; the per-direction
scaling $c = g^\top B_1 g / g^\top A_1 g$ restores a valid $\chi^2$ under the
local null. This is a \emph{denominator} correction.

\paragraph{Structural misspecification} (the model is wrong in direction $k \neq
j$). The score $s_j(\Th_0)$ is evaluated at the constrained solution, which the
estimator has biased by spreading the $k$-misfit across every frozen parameter.
This is a \emph{numerator} contamination, and no sandwich touches it.

\begin{proposition}[Why robustifying the MI cannot fix structural leakage]
Under a fixed structural error in direction $k$, the contaminating part of
$s_j(\Th_0)$ is $O(\sqrt{N})$ (its noncentrality grows with $N$; that is how the
test gains power against $k$), whereas any meat/bread reweighting is an $O(1)$
rescaling of the statistic. An $O(1)$ factor cannot recenter an $O(\sqrt{N})$
bias.
\end{proposition}

The robust MI is therefore correctly sized under the local null
(``$j$ constrained, rest correct, distribution possibly wrong'') and silent
about structural bias. The score test's anchor model is the constrained one,
which under genuine non-invariance \emph{is} the misspecified model, so its
per-item MIs inherit the leakage. The Wald's anchor (the free model) is, for a
loadings-only alternative, correctly specified; this is the structural reason a
Wald localization resists the contamination a score MI suffers.

\section{The LRT modification index}

To decontaminate the numerator one must stop conditioning on the suspect
parameters at wrong fixed values. The LRT does this by \emph{refitting}: for each
candidate constraint, fit the released model $H_1$ and report
\[
  T_d = T_0 - T_1, \qquad \mathrm{df} = m,
\]
together with the robust reference-law family from the Satorra (2000) difference
spectrum. Let $\lambda_1,\dots,\lambda_m$ be the eigenvalues of the difference
operator; then
\[
  T_d \dto \sum_{j=1}^m \lambda_j \chi^2_{1,j},
\]
and magmaan reports the naive $\chi^2_m$ tail, the Satorra--Bentler mean scaling
$\bar c = \tfrac1m\sum_j\lambda_j$, the mean--variance adjustment, the
scaled-shifted statistic, the exact mixture tail (Imhof/Farebrother), and the
penalized FMG/pEBA transform. None of these cost extra: they are functions of the
same spectrum the nested difference test already computes.

\paragraph{Exact EPC.} The one-step expected parameter change is a single Newton
step, $\mathrm{epc} = s_{\mathrm{eff}}/\mathcal{I}_{\mathrm{eff}}$. The LRT
counterpart is the \emph{refitted} value of the freed parameter minus its
constrained value: no linearization, the actual change. The existing
standardization (\texttt{epc.lv}, \texttt{epc.all}) applies unchanged.

\section{Conditional localization: exact, not local}

Because the LRT names both endpoints, it can \emph{condition}. To test item $j$
while protecting against a suspected non-invariant set $B$, free $B$ in
\emph{both} $H_0$ and $H_1$ (in magmaan, \texttt{group\_partial}). Then $B$'s
misfit is identical on the two sides and cancels in $T_d$, so $j$ is isolated.

This is the refit analogue of the Bera--Yoon (1993) adjusted score test, which
orthogonalizes $s_j$ against the $k$-direction. The crucial difference: Bera--Yoon
removes the \emph{linear} dependence and is valid only under local
($O(N^{-1/2})$) misspecification; the conditional LRT removes $B$'s misfit to all
orders by genuinely refitting, so it is exact under fixed non-invariance. The two
extremes of the conditioning dial are the ordinary MI (condition on everything,
maximal contamination) and the configural Wald (condition on nothing, zero
contamination, all df spent).

\begin{remark}
The residual cost is the specification-search/anchor-item problem: a third
non-invariant item left in the tied background contaminates both the $j$-test and
the choice of $B$. And the difference spectrum is evaluated at $H_1$'s pseudo-true
point; the exact mixture fixes the reference-law \emph{shape}, while conditioning
fixes the \emph{point} at which it is evaluated. These are orthogonal jobs.
\end{remark}

\section{Cost}

The refit sweep is $\#\text{candidates}$ nested fits plus one difference spectrum
each, against the one fit of the score sweep. Table~\ref{tab:cost} times both for
a two-group one-factor CFA (loadings $0.6$, one planted non-invariant loading,
$n_A:n_B = 1{:}3$), cold-started (magmaan computes its own start values).

\begin{table}[h]
\centering
\caption{Cold-start LRT-MI vs.\ one-step score sweep, $n=1000$, median of 3 runs.}
\label{tab:cost}
\begin{tabular}{rrrrrrr}
\toprule
items & npar & cand & one-step (ms) & LRT (ms) & per cand (ms) & LRT/one-step \\
\midrule
 6 & 24 &  5 & 1.0 &  6.0 & 1.2 & 6.0 \\
 9 & 36 &  8 & 1.0 & 15.0 & 1.9 & 15.0 \\
12 & 48 & 11 & 1.0 & 23.0 & 2.1 & 23.0 \\
15 & 60 & 14 & 2.0 & 43.0 & 3.1 & 21.5 \\
20 & 80 & 19 & 6.0 & 97.0 & 5.1 & 16.2 \\
\bottomrule
\end{tabular}
\end{table}

The ratio LRT/one-step is $6$--$23\times$, but only because the one-step sweep is
near-free; in absolute terms the entire LRT sweep is sub-100\,ms even at 20 items.
Sample size enters mildly through the empirical $\hat\Gamma$: at 12 items the
per-candidate cost grows from $2.0$\,ms at $n=500$ to $4.7$\,ms at $n=8000$. The
conclusion for the design question is unambiguous: at the model sizes
measurement-invariance studies use, the refit cost is not a barrier.

\paragraph{Implementation.} The sweep ships as the R function
\texttt{modification\_indices\_lrt(fit, data)}, a thin composition of the
one-call estimator \texttt{magmaan} and the nested test
\texttt{robust\_nested\_lrt}: it enumerates the tied families from the anchor
partable, refits each released model, and returns the difference statistic, the
reference-law family, and the exact released per-group estimates. An existing
\texttt{group\_partial} baseline on the anchor is carried into every refit, so
the conditional localization is automatic. For both complete-data ML and FIML the
per-row statistic reproduces a hand-built two-fit \texttt{robust\_nested\_lrt} to
machine precision.

\paragraph{What remains as an optimization.} This baseline cold-starts every
released fit. Starting each released fit from the constrained MLE (the C++ core
takes an explicit \texttt{x0} and contracts it into the released parameter space)
should cut the per-candidate refit to a few Newton steps. The benchmark
(\texttt{benchmarks/r/bench\_mi\_lrt.R}) is structured to quantify that warm-start
speed-up against this cold-start reference once the warm path lands.

\end{document}
